% Owner: Ankit (Member 3)
%=====================================================================
\section{Literature Review}
%=====================================================================
\subsection{Stored-program memory organisations}
In von Neumann's EDVAC design, instructions and data are stored in the same memory
\cite{vonneumann1945}. Programs can therefore be loaded and changed like any other data, but
every instruction needs at least one memory transfer (its fetch), and loads and stores need a
second one. Backus called the single connection between processor and store the
``von~Neumann bottleneck'' \cite{backus1978}, and because processor speed has grown much faster
than memory speed, its cost keeps rising, a trend described as the memory wall \cite{wulf1995}.
The simplest visible form of the bottleneck is a structural hazard: when a fetch and a data
access need the same port in the same cycle, one of them must wait \cite{hp_quant}.

The Harvard organisation takes its name from the Harvard Mark~I, whose instructions came from
punched tape while numbers were held in separate counters \cite{aiken1946}. Physically separate
instruction and data memories, each with its own address and data bus, double the bandwidth
available to the processor, but the division of storage between code and data is fixed and a
second set of buses is needed \cite{hp_quant}. Textbook single-cycle processors need exactly this
arrangement, because a load must read an instruction word and a data word in the same clock
period, whereas a multi-cycle processor can share one memory by fetching and accessing data in
different cycles \cite{ph_riscv,harris_riscv}. Our Phase~I designs are these two textbook cases.

\subsection{Modified Harvard organisations}
General-purpose processors rarely keep two separate main memories. They apply the Harvard idea
only at the top of the hierarchy: split level-1 instruction and data caches give the pipeline
two ports, and a single unified memory below the caches holds both code and data, so programs can
still be loaded as data \cite{hp_quant,ph_riscv}. Smith's survey already compares split and
unified caches. The split gives twice the bandwidth and lets each cache be optimised for its
access pattern. The unified cache shares its capacity dynamically between code and data, so a
program whose code is large and data small, or the reverse, can still use all of it
\cite{smith1982}. Which effect wins depends on the workload, on the cache size and on whether the
processor can actually use two accesses per cycle. Phase~II measures this trade-off directly.
